\documentclass[11pt,letterpaper]{article}
\usepackage[margin=1in]{geometry}
\usepackage{iftex}
\ifPDFTeX
  \usepackage[T1]{fontenc}
  \usepackage[lining,tabular]{carlito}
  \renewcommand{\familydefault}{\sfdefault}
  \newcommand{\headingfont}{\fontfamily{ptm}\selectfont}
\else
  \usepackage{fontspec}
  \setmainfont{Carlito}
  \newfontfamily\headingfont{Times New Roman}
\fi
\usepackage{graphicx}
\usepackage{xcolor}
\usepackage{eso-pic}
\pagestyle{empty}
\frenchspacing
\setlength{\parindent}{0pt}
\setlength{\parskip}{0pt}
\setlength{\unitlength}{1bp}
\newcommand{\Text}[3]{\put(#1,#2){\makebox(0,0)[lb]{\smash{#3}}}}
\newcommand{\Justified}[4]{\Text{#1}{#2}{{\spaceskip=\fontdimen2\font plus 1fil minus 1pt\relax\makebox[#3bp][s]{#4}}}}
\newcommand{\Rule}[4]{\put(#1,#2){\rule{#3bp}{#4bp}}}
\newcommand{\ResponseFrame}[2]{%
  \Rule{71.76}{#1}{0.48}{\fpeval{#2-#1}}%
  \Rule{539.52}{#1}{0.48}{\fpeval{#2-#1}}%
  \Rule{71.76}{#1}{468.24}{0.48}%
  \Rule{71.76}{\fpeval{#2-0.96}}{468.72}{0.96}}
\newcommand{\Divider}[1]{\Rule{71.76}{#1}{468.24}{0.48}}
\newcommand{\Bullet}[1]{\Text{90}{#1}{\textbullet}}
\newcommand{\FormPage}[2]{%
  \null
  \AddToShipoutPictureFG*{\AtPageLowerLeft{%
    \setlength{\unitlength}{1bp}%
    \fontsize{10.98bp}{14.5bp}\selectfont
    \begin{picture}(0,0)
      #2
      \Text{303.2}{51.72}{#1}
    \end{picture}}}
  \newpage}

\begin{document}

\FormPage{1}{
  \put(79.9,620.82){\includegraphics[width=97.6bp]{prism-uploads/uctLogo (1).png}}
  \Rule{196.28}{620.8}{0.96}{99.2}
  \Rule{71.28}{611.26}{468.54}{0.96}
  \Text{213.8}{705.06}{{\headingfont\fontsize{16.02bp}{26.4bp}\selectfont UNIVERSITY OF CAPE TOWN}}
  \Text{213.8}{678.66}{{\headingfont\fontsize{16.02bp}{26.4bp}\selectfont Department of Electrical Engineering}}
  \Text{213.8}{652.24}{{\headingfont\fontsize{16.02bp}{26.4bp}\selectfont EEE4022F/S - Final Year Project}}
  \Text{213.8}{625.84}{{\headingfont\fontsize{16.02bp}{26.4bp}\selectfont Graduate Attribute Tracking Form}}

  \Text{77.4}{576.52}{Student name:}
  \Text{170}{576.52}{Samson Okuthe}
  \Text{312.62}{576.52}{DP Awarded? [Y/N]}
  \Text{77.4}{558.52}{Student no:}
  \Text{170}{558.52}{OKTSAM001}
  \Text{312.62}{558.52}{Supervisor name:}
  \Text{417}{558.52}{Simon Winberg}
  \Text{77.4}{540.52}{Date:}
  \Text{170}{540.52}{\today}
  \Text{312.62}{540.52}{Date:}
  \Text{417}{540.52}{\today}
  \Text{77.4}{509.74}{Student signature:}
  \Text{312.62}{516.46}{Supervisor}
  \Text{312.62}{503.08}{signature:}
  \Rule{165.5}{570.82}{127.56}{0.96}
  \Rule{412.12}{570.82}{127.58}{0.96}
  \Rule{165.5}{552.82}{127.56}{0.96}
  \Rule{412.12}{552.82}{127.58}{0.96}
  \Rule{165.5}{534.82}{127.56}{0.96}
  \Rule{412.12}{534.82}{127.58}{0.96}
  \Rule{164.78}{491.26}{128.28}{0.96}
  \Rule{411.4}{491.26}{128.3}{0.96}

  \Text{72}{458.32}{\textbf{Instructions:}}
  \Justified{72}{435.82}{468.011}{Students must explain in this document what they \textbf{have already done} and what they \textbf{plan to do} to satisfy}
  \Justified{72}{421.3}{467.912}{each Graduate Attribute. Descriptions of each GA is provided at the end of the document. Supervisors}
  \Justified{72}{406.84}{467.748}{respond to the student\textquotesingle s plans and current progress, providing additional comments or advice as they see}
  \Justified{72}{392.32}{467.506}{fit. Once the student\textquotesingle s progress is deemed sufficient (a few weeks before submission at the due date for}
  \Text{72}{377.84}{this form), supervisors indicate that DP can be awarded.}
  \Text{72}{355.34}{\textcolor{red}{\textbf{VERY IMPORTANT: Receiving DP for the course does NOT imply that all GA’s have been met in the}}}
  \Text{72}{340.82}{\textcolor{red}{\textbf{course. Assessment of GA’s only happen in the final marking of the project report.}}}

  \Text{255.2}{318.38}{\textbf{GA 1: Problem Solving}}
  \ResponseFrame{83.04}{311.84}
  \Divider{168.62}
  \Text{77.4}{300.38}{Student Response:}
  \put(77.4,293){\parbox[t]{462bp}{\fontsize{8.3bp}{10bp}\selectfont
  I formulated the core research question around measurable vehicle dynamic response rather than subjective ride comfort, rejecting a driver-comfort framing in favour of one suited to signal-processing analysis. I root-caused several faults from first principles: a GPIO conflict I resolved via the manufacturer\textquotesingle s own schematic; a BLE service-discovery failure I confirmed against an independent reference implementation; and a systematic measurement flaw I caught before any field data existed to be corrupted by it, where my literature review showed that a raw acceleration RMS conflates road roughness with vehicle speed, since the road-to-acceleration transfer function is itself speed dependent. I then conducted a competitive-landscape and literature audit to formulate a specific, evidenced research gap: synchronised multi-sensor vehicle-dynamics characterisation for a three-wheeled cargo vehicle, distinct from existing four-wheel and motorcycle telemetry work and existing commercial telematics-plus-HUD products. The next stage is to run this characterisation on the moving vehicle rather than the bench, since every subsystem validated so far has been proven individually but not yet on the vehicle the research question is actually about.}}
  \Text{77.4}{158.12}{Supervisor Response:}
}

\FormPage{2}{
  \Text{185.84}{709.5}{\textbf{GA 4: Investigations, Experiments, and Data Analysis}}
  \ResponseFrame{474.16}{703.02}
  \Divider{559.72}
  \Text{77.4}{691.5}{Student Response:}
  \put(77.4,684){\parbox[t]{462bp}{\fontsize{8.3bp}{10bp}\selectfont
  I conducted hypothesis-test-conclusion investigations throughout: the GNSS field-count assumption, which I hypothesised at 13 fields, tested live, found to be 9, and corrected; the BMS service-UUID assumption, which I tested against a powered pack, found wrong, corrected, and cross-validated against an independent parser; and the sync-pulse timing offset, which I bounded to under one second via a lens-cover test. I designed a controlled experiment for the ride-characterisation feature, using a synthetic sine wave with analytically-known RMS, peak-to-peak, and crest-factor values, validating the computation independently of sensor hardware. A further investigation of mine found two nominally-identical accelerometers disagreeing by about 7\% at rest; I hypothesised this as factory tolerance, corrected it via a per-sensor calibration, and confirmed it held under an induced dynamic disturbance too, not only the static case it was derived from. My continuously-run literature pipeline has retrieved 733 candidate sources and confirmed the DOI of 114 of them against the published title. The next stage is to bench-validate the newly sourced VL53L5CX proximity sensor against an independent published characterisation of the same chip before trusting its datasheet figures, and to capture a genuine outdoor GNSS fix to resolve the one remaining untested field-mapping assumption.}}
  \Text{77.4}{549.22}{Supervisor Response:}

  \Text{236.12}{441.22}{\textbf{GA 5: Use of Engineering Tools}}
  \ResponseFrame{205.88}{434.68}
  \Divider{291.44}
  \Text{77.4}{423.22}{Student Response:}
  \put(77.4,416){\parbox[t]{462bp}{\fontsize{8.3bp}{10bp}\selectfont
  I applied, and recognised the limitations of, PlatformIO, Arduino, and FreeRTOS for embedded development; \texttt{esptool} and serial diagnostics for hardware verification; OpenSCAD for a parametric, printable camera-mount design; and Python automation (\texttt{paramiko}, \texttt{PyMuPDF} for schematic verification, \texttt{PIL} for frame analysis) that I built to support testing. I verified tool limitations explicitly rather than assuming them: I did not accept a BMS checksum formula as correct merely because a library accepted it, with no vendor datasheet available to confirm it independently; and I re-checked an OpenSCAD export metric, initially misread as a mount-design fault, against a known-valid control case (a plain reference cube) before accepting or rejecting the concern. The next stage is to complete the LittleFS replay-on-reconnect logic and integrate the VL53L5CX proximity sensor into the existing per-task acquisition architecture once it arrives.}}
  \Text{77.4}{280.94}{Supervisor Response:}

  \Text{194.18}{172.88}{\textbf{GA 6: Professional and Technical Communication}}
  \ResponseFrame{108.2}{166.4}
  \Text{77.4}{154.88}{Student Response:}
  \put(77.4,148){\parbox[t]{462bp}{\fontsize{7.6bp}{9bp}\selectfont
  I have maintained a continuously-updated Methodology chapter documenting my engineering work honestly, including faults found and how I resolved them, and negotiated shared vehicle access in writing with the Controls Lab custodian. I also caught and corrected a cross-chapter inconsistency before it reached a reader: the Introduction described the camera's object detector as an active feature, while my Methodology had documented it as withheld under the ethics constraint. I am now drafting the Results chapter directly from the evidence already logged in Methodology.}}
}

\FormPage{3}{
  \Rule{71.76}{676.5}{0.48}{43.5}
  \Rule{539.52}{676.5}{0.48}{43.5}
  \Divider{719.52}
  \Divider{676.5}
  \Text{77.4}{709.02}{Supervisor Response:}

  \Text{248.9}{643.48}{\textbf{GA 8: Individual Working}}
  \ResponseFrame{578.32}{637}
  \Rule{71.52}{534.76}{0.96}{44.52}
  \Rule{539.28}{534.76}{0.96}{44.52}
  \Rule{71.52}{534.76}{468.72}{0.96}
  \Rule{71.52}{578.32}{468.72}{0.96}
  \Text{77.4}{625.48}{Student Response:}
  \put(77.4,619){\parbox[t]{462bp}{\fontsize{7.6bp}{9bp}\selectfont
  I took full individual responsibility for hardware bring-up, firmware debugging, and report authorship, making technical decisions independently and bringing my supervisor specific questions rather than requests for direction. When a nearby-object alert for the driver became a requirement, I chose an 8$\times$8-zone time-of-flight sensor over a single-point one to give left, centre and right discrimination, and withheld the camera detector until its ethics approval came through. I respected shared resources, leaving the inherited dashboard\textquotesingle s SD card and another team\textquotesingle s vehicle-access slots untouched. I am now coordinating joint integration onto the vehicle ahead of the field test.}}
  \Text{77.4}{567.82}{Supervisor Response:}

  \Text{225.86}{501.82}{\textbf{GA 9: Independent Learning Ability}}
  \ResponseFrame{294.8}{495.28}
  \Divider{380.42}
  \Text{77.4}{483.82}{Student Response:}
  \put(77.4,476){\parbox[t]{462bp}{\fontsize{8.3bp}{10bp}\selectfont
  I repeatedly corrected my own assumptions against independently-gathered evidence: I verified GPIO pin assignments against a primary-source schematic after a placeholder proved wrong, and I corrected a vehicle-identity error after independently checking documentation. I built a continuously-running literature pipeline specifically to keep self-correcting rather than trust one search; it has since surfaced real prior art that forced two of my research-gap claims to narrow rather than stand unexamined. I also checked the modem's behaviour against the manufacturer's 652-page AT command manual; this showed that the documented GNSS response has sixteen fields and reports speed in knots, while the fitted modem returns nine, so I now treat the manual as a hypothesis to test against the hardware rather than as ground truth. The next stage is to apply the same evidence-before-assumption standard to the newly adopted proximity sensor, treating its independently published characterisation as the basis for expected performance rather than its datasheet figures alone.}}
  \Text{77.4}{369.92}{Supervisor Response:}

  \Text{72}{261.86}{\textbf{GA 1: Problem Solving}}
  \Justified{72}{239.36}{409.203}{Identify, formulate, research literature and analyse complex engineering problems reaching}
  \Justified{72}{224.84}{431.749}{substantiated conclusions using first principles of mathematics, natural sciences and engineering}
  \Justified{72}{210.38}{295.087}{sciences with holistic considerations for sustainable development.}
  \Bullet{195.32}
  \Justified{108.02}{195.32}{413.655}{A systematic, theory-based understanding of the natural sciences applicable to the discipline}
  \Text{108.02}{180.8}{and awareness of relevant social sciences.}
  \Bullet{165.74}
  \Justified{108.02}{165.74}{429.076}{Conceptually based mathematics, numerical analysis, data analysis, statistics and formal aspects}
  \Justified{108.02}{151.22}{417.482}{of computer and information science to support detailed analysis and modelling applicable to}
  \Text{108.02}{136.76}{the discipline.}
  \Bullet{121.7}
  \Justified{108.02}{121.7}{377.300}{A systematic, theory-based formulation of engineering fundamentals required in the}
  \Text{108.02}{107.18}{engineering discipline.}
}

\FormPage{4}{
  \Bullet{708.96}
  \Justified{108.02}{708.96}{431.075}{Engineering specialist knowledge that provides theoretical frameworks and bodies of knowledge}
  \Justified{108.02}{694.44}{430.592}{for the accepted practice areas in the engineering discipline, much of which is at the forefront of}
  \Text{108.02}{679.92}{the discipline.}

  \Text{72}{650.92}{\textbf{GA 4: Investigations, Experiments and Data Analysis}}
  \Justified{72}{628.42}{467.952}{Demonstrate competence to conduct investigations of complex engineering problems using research}
  \Justified{72}{613.96}{467.904}{methods, including research-based knowledge, design of experiments, analysis and interpretation of data,}
  \Justified{72}{599.44}{256.971}{and synthesis of information to provide valid conclusions.}
  \Justified{72}{570.46}{467.910}{Engagement with selected knowledge in the current research literature of the discipline, awareness of}
  \Justified{72}{556}{370.343}{the power of critical thinking and creative approaches to evaluate emerging issues.}
  \Justified{72}{527.02}{467.660}{The balance of investigation and experiment should be appropriate to the discipline. Research}
  \Justified{72}{512.5}{467.671}{methodology to be applied in research or investigation where the student engages with selected}
  \Text{72}{498.04}{knowledge in the research literature of the discipline.}
  \Justified{72}{469.06}{467.660}{Note: An investigation differs from a design in that the objective is to produce knowledge and}
  \Justified{72}{454.54}{467.594}{understanding of a phenomenon and a recommended course of action rather than specifying how an}
  \Text{72}{440.08}{artefact could be produced.}

  \Text{72}{403.06}{\textbf{GA 5: Use of engineering tools}}
  \Justified{72}{380.54}{418.865}{Demonstrate competence to create, select and apply and recognise limitations of appropriate}
  \Justified{72}{366.08}{440.107}{techniques, resources and modern engineering and IT tools, including prediction and modelling, to}
  \Text{72}{351.56}{complex engineering problems.}
  \Bullet{336.56}
  \Justified{108.02}{336.56}{429.076}{Conceptually based mathematics, numerical analysis, data analysis, statistics and formal aspects}
  \Justified{108.02}{322.04}{417.482}{of computer and information science to support detailed analysis and modelling applicable to}
  \Text{108.02}{307.52}{the discipline.}
  \Bullet{292.52}
  \Justified{108.02}{292.52}{390.515}{Knowledge of engineering practice (technology) in the practice areas in the engineering}
  \Text{108.02}{277.94}{discipline}
  \Justified{72}{263.48}{451.492}{A range of techniques, resources and modern engineering and IT tools appropriate to the disciplinary}
  \Text{72}{248.96}{designation of the programme.}

  \Text{72}{212}{\textbf{GA 6: Professional and Technical Communication}}
  \Justified{72}{189.5}{467.671}{Demonstrate competence to communicate effectively, both orally and in writing, with engineering}
  \Justified{72}{175.04}{460.607}{audiences and the community at large, taking into account cultural, language, and learning differences.}
  \Justified{72}{152.54}{467.572}{This course evaluates the long report component of this outcome at exit level.\ \ Material to be}
  \Justified{72}{138.02}{467.572}{communicated is in an academic or simulated professional context. Audiences range from engineering}
  \Justified{72}{123.56}{467.925}{peers, management and lay persons, using appropriate academic or professional discourse. Written}
  \Justified{72}{109.04}{467.965}{reports (10 000 to 15 000 words plus tables, diagrams and appendices) should cover material at exit-level.}
  \Justified{72}{94.56}{467.638}{Methods of providing information include the conventional methods of the discipline, for example}
  \Justified{72}{80.04}{259.948}{engineering drawings, as well as subject-specific methods.}
}

\FormPage{5}{
  \Text{72}{687}{\textbf{GA 8: Individual, Team and Multidisciplinary Working}}
  \Justified{72}{664.5}{467.660}{Demonstrate competence to work effectively as an individual, in teams and in multidisciplinary}
  \Justified{72}{650.02}{467.851}{environments. This course evaluates the \textbf{individual} working component of this learning outcome at exit}
  \Text{72}{635.5}{level.}
  \Justified{72}{613}{378.085}{Knowledge of professional ethics, responsibilities and norms of engineering practice.}

  \Text{72}{584.02}{\textbf{GA 9: Independent Learning Ability}}
  \Justified{72}{561.58}{446.886}{Demonstrate competence to engage in independent learning through well developed learning skills.}
  \Justified{72}{532.6}{467.910}{Engagement with selected knowledge in the current research literature of the discipline, awareness of}
  \Justified{72}{518.08}{370.343}{the power of critical thinking and creative approaches to evaluate emerging issues.}
  \Justified{72}{489.1}{467.723}{Operate independently in complex, ill-defined contexts requiring personal responsibility and initiative,}
  \Justified{72}{474.58}{467.662}{accurately self-evaluate and take responsibility for learning requirements; be aware of social and ethical}
  \Justified{72}{460.12}{257.207}{implications of applying knowledge in particular contexts.}
  \Bullet{430.6}
  \Justified{108.02}{430.6}{431.591}{Openness to constructive feedback, awareness of own limitations, ability to cope with the}
  \Justified{108.02}{416.08}{431.646}{discomfort of uncertainty and having access to a range of approaches, reflective selfevaluation,}
  \Justified{108.02}{401.56}{431.959}{curiosity and proactive engagement, resilience, confidence to ask for help and draw from a broad}
  \Text{108.02}{387.1}{range of stakeholders.}
  \Bullet{372.02}
  \Justified{108.02}{372.02}{431.784}{Reflection of self-learning to begin to recognise if what has been covered meets the needs of the}
  \Text{108.02}{357.5}{activity or task.}
}

\end{document}
